\documentclass[11pt]{article}
\usepackage[margin=1in]{geometry}
\usepackage{amsmath, amssymb}
\usepackage{enumitem}
\usepackage{titlesec}
\usepackage{hyperref}
\usepackage{booktabs}
\usepackage{xcolor}
\usepackage{parskip}

\hypersetup{colorlinks=true, linkcolor=blue, urlcolor=blue}

\title{\textbf{EE627 --- Class Summary}\\[0.4em]
       \large January 30, 2026}
\author{Prof.\ Rensheng Wang}
\date{}

\begin{document}
\maketitle
\tableofcontents
\newpage

%------------------------------------------------------------
\section{XGBoost: What It Is and Why It Matters}
%------------------------------------------------------------

\subsection{Definition}
\textbf{XGBoost} (eXtreme Gradient Boosting) is a decision-tree-based \emph{ensemble} machine
learning algorithm that uses a gradient boosting framework.

\begin{itemize}
    \item Open-source; core implemented in C++, with R/Python/Julia interfaces.
    \item Cloud-ready: supports AWS, Azure, Yarn; integrates with Flink and Spark.
    \item \textbf{The winning model for numerous Kaggle competitions.}
\end{itemize}

\subsection{When to Use It}
\begin{itemize}
    \item \textbf{Structured (tabular) data} --- CSV, TSV, SQL tables with named columns.
      XGBoost is the \emph{first model to try} when you have no strong prior.
    \item \textbf{Not for unstructured data} (images, text). Those require deep learning
      (neural networks, transformers) because pixel or token positions carry no fixed
      semantic meaning that a tree can split on.
\end{itemize}

\subsection{Why This Is an Interview Topic}
XGBoost is a high-frequency data engineering interview question because:
\begin{enumerate}
    \item It is not a yes/no question --- it invites extended discussion.
    \item Explaining \emph{why} it works shows system-level understanding, not just API familiarity.
    \item Knowing its limits (when \emph{not} to use it) signals senior-level thinking.
\end{enumerate}

\newpage
%------------------------------------------------------------
\section{The Hiring Manager Analogy: Decision Tree to XGBoost}
%------------------------------------------------------------

The evolution of tree-based algorithms can be understood entirely through the lens of
an iterative interview process.

\subsection{Decision Tree --- One Interviewer}
A single hiring manager applies their own set of criteria: education level, years of
experience, coding ability, communication. Each criterion is a \emph{split} in the tree.
Easy to visualize and explain (interpretable), but \textbf{biased} toward one person's
perspective.

\subsection{Bagging (Bootstrap Aggregating) --- Interview Panel}
Replace the single interviewer with a \textbf{committee}. Each member votes independently.
The final decision aggregates all votes democratically.
\begin{itemize}
    \item Reduces individual bias.
    \item Majority vote $\Rightarrow$ more robust, statistically reliable outcome.
    \item If 9 of 10 panelists approve: strong positive signal. Mixed votes: doubt.
\end{itemize}

\subsection{Random Forest --- Randomized Panel}
Bagging, but each panelist only tests the candidate on a \emph{randomly chosen subset}
of criteria (e.g., one tests coding, another tests system design, another tests
behavioral). This randomness:
\begin{itemize}
    \item Prevents any single feature from dominating.
    \item Helps the ensemble generalize to new data (avoids overfitting).
    \item A few bad random subsets are harmless when many subsets are sampled.
\end{itemize}

\subsection{Boosting --- Sequential Interviews with Feedback}
Interviewers run sequentially. Before each session, the next interviewer receives
\textbf{feedback from the previous one} and sharpens their questions accordingly.
Each round focuses on the \emph{residual gap} --- the parts the previous round did
not resolve.
\begin{itemize}
    \item Feedback loop $\Rightarrow$ progressively more targeted evaluation.
    \item The model ``learns from its mistakes'' round by round.
\end{itemize}

\subsection{Gradient Boosting --- Measure the Gap with Calculus}
Boosting, but each round measures the residual gap using the \textbf{gradient} of the
loss function --- i.e., how far the current prediction is from the target, and in
which direction to push.

\textit{Analogy from lecture:} A candidate is great at communication but terrible at
coding. The gradient tells you exactly \emph{how large} that coding gap is. The next
interview round focuses precisely there.

\subsection{XGBoost --- Gradient Boosting on Steroids}
XGBoost improves on gradient boosting by also measuring the \textbf{Hessian} --- the
second derivative of the loss --- telling the model \emph{how quickly the gradient
itself is changing}.

\begin{itemize}
    \item First derivative (gradient $g_i$): how far off the prediction is.
    \item Second derivative (Hessian $h_i$): how fast that error is changing --- how
      aggressively to step toward the minimum.
    \item Combined, these enable fewer iterations and faster convergence than standard
      gradient boosting.
\end{itemize}

This is combined with state-of-the-art software and hardware optimizations (see
Section~\ref{sec:why_xgb_performs}).

\newpage
%------------------------------------------------------------
\section{Why XGBoost Performs So Well}
\label{sec:why_xgb_performs}
%------------------------------------------------------------

XGBoost improves on standard Gradient Boosting Machines (GBMs) through
\textbf{system optimizations} and \textbf{algorithmic enhancements}.

\subsection{System Optimizations}

\subsubsection{Parallelization}
Standard tree building is sequential (outer loop over leaf nodes, inner loop over
features). XGBoost swaps this loop order, enabling features to be evaluated
simultaneously across parallel threads.

\textit{Why it matters:} From 100 features, choosing a random subset of 20 yields
a combination space larger than $10^{15}$. Sequential evaluation would be
computationally infeasible. Parallelization makes it tractable.

\subsubsection{Tree Pruning (Depth-First)}
Standard GBMs use a greedy stopping criterion: stop splitting when the split would
be negative. This is ``short-sighted'' --- a seemingly bad split now might enable a
much better split one level deeper.

XGBoost uses \texttt{max\_depth} and grows trees \textbf{depth-first}, then
\emph{prunes backward}. This finds globally better tree structures.

\subsubsection{Hardware Optimization (Cache Awareness)}
Each thread allocates an internal buffer to store gradient statistics locally. This
avoids redundant memory accesses --- values computed once are cached and reused.

\textit{Connection:} The same cache-awareness principle drives the key-value (KV)
cache in Transformer models, where intermediate attention computations are stored
and reused rather than recalculated.

\subsection{Algorithmic Enhancements}

\subsubsection{Regularization (Built-In L1 + L2)}
XGBoost's objective penalizes model complexity with \emph{both} L1 (Lasso) and L2
(Ridge) regularization simultaneously:
\[
\mathcal{L}(\phi) = \sum_i l(\hat{y}_i,\, y_i) + \sum_k \Omega(f_k),
\quad \Omega(f) = \gamma T + \tfrac{1}{2}\lambda\|w\|^2
\]
where $T$ is the number of leaves and $w$ are the leaf weights. You do \emph{not}
need to add a regularizer manually --- it is baked in.

\textbf{L1 vs.\ L2 (a frequent interview question):}
\begin{itemize}
    \item \textbf{L1 (Lasso, \texttt{alpha}):} penalizes $|w_j|$.
      Geometrically, the constraint region has corners, which forces many weights
      to exactly zero $\Rightarrow$ \textbf{fast, aggressive feature elimination}.
      Use when you want dimension reduction.
    \item \textbf{L2 (Ridge, \texttt{lambda}):} penalizes $w_j^2$.
      Constraint region is smooth; weights shrink toward zero but rarely hit it
      exactly $\Rightarrow$ \textbf{slower, gentler shrinkage}. Keeps all features.
\end{itemize}

\subsubsection{Sparsity Awareness (Missing Value Handling)}
When a value is missing in the feature matrix, XGBoost learns the optimal
\emph{default direction} (left or right branch) from the training data itself.
The algorithm tries both directions and keeps whichever yields better performance.

\textbf{Practical result:} You do \emph{not} need to impute missing values before
running XGBoost. The model handles them natively and intelligently.

\textit{Contrast with legacy tools:} MATLAB and older Python libraries returned
errors on missing data, forcing analysts to delete rows or manually impute values
(mean, zero, or custom). Both choices introduce bias. XGBoost eliminates this pain.

\subsubsection{Weighted Quantile Sketch}
Before building any tree, XGBoost pre-scans the data to compute \emph{percentile
breakpoints} (quantiles) for each feature. Tree splits are proposed only at these
breakpoints, not at arbitrary values.

\textit{Why:} A split at age $= 10$ is meaningless if everyone in the dataset is
over 18. By understanding the data distribution first, every proposed split is
statistically defensible and the same quantile information can be reused across
multiple trees without recomputation.

\subsubsection{Built-In Cross-Validation}
XGBoost performs internal train/validation splitting at each boosting iteration.
You do not need to manually code a 70/30 split or $k$-fold cross-validation ---
it is already happening inside the model.

\newpage
%------------------------------------------------------------
\section{Benchmark: Where Is the Proof?}
%------------------------------------------------------------

Using scikit-learn's \texttt{Make\_Classification} dataset (1 million points, 20
features, 5-fold CV):

\begin{center}
\begin{tabular}{@{}lrr@{}}
\toprule
\textbf{Algorithm}   & \textbf{AUROC} & \textbf{Training Time (s)} \\
\midrule
XGBoost              & 0.9662         & \textbf{24}               \\
Gradient Boosting    & 0.9661         & 2{,}069                   \\
Random Forest        & 0.9542         & 424                       \\
Logistic Regression  & 0.9373         & 17                        \\
\bottomrule
\end{tabular}
\end{center}

XGBoost matches the best accuracy while being \textbf{86$\times$ faster} than
standard Gradient Boosting.

\newpage
%------------------------------------------------------------
\section{The ``No Free Lunch'' Caveat}
%------------------------------------------------------------

XGBoost is the default starting point, not a guaranteed winner. Real practice
requires:

\begin{enumerate}
    \item Testing multiple algorithms on your specific data.
    \item \textbf{Hyperparameter tuning} --- a model with default settings is never
      fully optimized.
    \item Considering computational complexity, explainability requirements, and
      ease of deployment.
\end{enumerate}

\textit{From the slides:} ``This is exactly the point where Machine Learning starts
drifting away from science towards art --- but honestly, that's where the magic
happens.''

\subsection{Key Hyperparameters to Know}

\begin{center}
\begin{tabular}{@{}lll@{}}
\toprule
\textbf{Parameter}       & \textbf{Controls}                              & \textbf{Typical Range} \\
\midrule
\texttt{learning\_rate}  & Step size per boosting round                   & 0.01 -- 0.3            \\
\texttt{n\_estimators}   & Number of boosting rounds (trees)              & 100 -- 1{,}000         \\
\texttt{max\_depth}      & Maximum depth of each tree                     & 3 -- 10                \\
\texttt{min\_child\_weight} & Min sum of instance weight to split further & 1 -- 10                \\
\texttt{subsample}       & Fraction of rows sampled per tree              & 0.5 -- 1.0             \\
\texttt{colsample\_bytree} & Fraction of columns sampled per tree         & 0.5 -- 1.0             \\
\texttt{lambda}          & L2 regularization weight                       & 0 -- 10                \\
\texttt{alpha}           & L1 regularization weight (Lasso)               & 0 -- 1                 \\
\texttt{gamma}           & Min loss reduction required to make a split    & 0 -- 5                 \\
\bottomrule
\end{tabular}
\end{center}

\textit{Practical advice from lecture:} Most parameters will show little sensitivity
on a given dataset. The value of tuning is to \emph{find the one or two that are
sensitive} --- those are the parameters that reveal the structure of your specific
data.

\subsection{Python Code}

\begin{verbatim}
# Regression
from xgboost import XGBRegressor
from sklearn.model_selection import cross_val_score
scores = cross_val_score(
    XGBRegressor(objective='reg:squarederror'),
    X, y, scoring='neg_mean_squared_error')

# Classification
from xgboost import XGBClassifier
cross_val_score(XGBClassifier(), X, y)
\end{verbatim}

\newpage
%------------------------------------------------------------
\section{XGBoost vs.\ LightGBM}
%------------------------------------------------------------

\subsection{What Is LightGBM?}
Microsoft's \textbf{LightGBM} (Light Gradient Boosting Machine) is a more efficient
variant of gradient boosting, now preferred by many Kagglers.

\begin{itemize}
    \item Up to $\mathbf{6\times}$ faster than XGBoost.
    \item Lower memory usage.
    \item Supports parallel and GPU learning.
    \item Capable of handling very large-scale data.
\end{itemize}

\subsection{The Core Difference: Leaf-Wise vs.\ Level-Wise Tree Growth}

\begin{center}
\begin{tabular}{@{}lll@{}}
\toprule
\textbf{Property}       & \textbf{XGBoost (Level-Wise)} & \textbf{LightGBM (Leaf-Wise)} \\
\midrule
Growth strategy         & Scan level by level (BFS)     & Always expand highest-loss leaf (DFS) \\
Speed                   & Slower                        & Faster                                 \\
Accuracy (small trees)  & Comparable                    & Often better                           \\
Accuracy (full tree)    & Same (converge to same tree)  & Same                                   \\
Monotonic constraints   & Supported                     & Not supported                          \\
Interpretability        & Easier                        & Slightly harder                        \\
\bottomrule
\end{tabular}
\end{center}

\textit{Key insight:} Leaf-wise growth chooses splits based on their contribution to
\emph{global} loss, not just the loss along one branch. For a fixed node budget, it
often produces lower-error trees than level-wise growth.

\subsection{LightGBM Novel Techniques}
\begin{itemize}
    \item \textbf{GOSS (Gradient-based One-Side Sampling):} Drops data instances
      with small gradients (they are already well-predicted). Keeps large-gradient
      instances and a random sample of small-gradient ones. Accurate information
      gain estimate at a fraction of the data size.
    \item \textbf{EFB (Exclusive Feature Bundling):} Bundles mutually exclusive
      features (features that are rarely non-zero simultaneously) into a single
      feature. Reduces the effective number of features, speeding up split finding.
\end{itemize}

\subsection{LightGBM Python Code}

\begin{verbatim}
from sklearn.model_selection import train_test_split
import lightgbm as lgb

X_train, X_test, y_train, y_test = \
    train_test_split(X, y, test_size=0.3, random_state=0)

clf = lgb.LGBMClassifier()
clf.fit(X_train, y_train)
\end{verbatim}

\subsection{Current Status}
XGBoost developers later added \texttt{grow\_policy='lossguide'} to match
LightGBM's leaf-wise behavior. Even so, LightGBM remains approximately
$1.3\times$--$1.5\times$ faster than XGBoost. Both are viable; XGBoost's
monotonic constraint support can be valuable when model interpretability is a
business requirement.

\newpage
%------------------------------------------------------------
\section{Correlation: The Foundation of Time Series}
\label{sec:correlation}
%------------------------------------------------------------

\subsection{Why Correlation?}
Correlation is the most fundamental measurement used throughout time series analysis.
Its applications are:
\begin{itemize}
    \item \textbf{Prediction:} if $X$ and $Y$ are correlated, knowing one lets you
      predict the other.
    \item \textbf{Feature selection:} highly correlated features are redundant;
      low-correlation features with the target are uninformative.
    \item \textbf{Reliability / validity:} measuring whether two instruments or
      raters agree.
\end{itemize}

\subsection{Pearson Correlation Coefficient}

The correlation coefficient between $X$ and $Y$ is:
\[
r_{x,y} = \frac{\text{Cov}(X,Y)}{\sqrt{\text{Var}(X)\,\text{Var}(Y)}}
         = \frac{E[(X-\mu_x)(Y-\mu_y)]}{\sqrt{E(X-\mu_x)^2\,E(Y-\mu_y)^2}}
\]

\begin{itemize}
    \item $-1 \leq r_{x,y} \leq 1$ always.
    \item $r \approx +1$: strong positive linear relationship.
    \item $r \approx -1$: strong negative linear relationship.
    \item $r \approx 0$: little or no \emph{linear} association.
\end{itemize}

\textbf{Critical subtlety:} $r = 0$ means \emph{uncorrelated}, \textbf{not}
\emph{independent}. Two variables can have a strong non-linear relationship
(e.g., $Y = X^2$) and still have $r = 0$.

\subsection{Computing Correlation in Python / MATLAB / R}

\begin{verbatim}
# Python (NumPy)
import numpy as np
r = np.corrcoef(x, y)[0, 1]

# MATLAB
r = corrcoef(A)          % correlation matrix of all columns

% R
cor(x, y)                # scalar; cor(A) gives full matrix
\end{verbatim}

\newpage
%------------------------------------------------------------
\section{Time Series Analysis: Stationarity}
%------------------------------------------------------------

\subsection{What Is a Time Series?}
A time series is data indexed by a sequential, \textbf{non-swappable} time index.
Order matters: shuffling destroys the dependency structure.

The same logic applies to any ordered sequence:
\begin{itemize}
    \item Financial prices (daily/hourly returns)
    \item Text tokens (word order in a sentence)
    \item Geospatial measurements (neighboring locations)
    \item Audio/spectral data (frequency bins)
\end{itemize}
This is why understanding time series transfers directly to understanding
language models and diffusion models.

\subsection{Stationarity}
The foundation of time series analysis. We need the process to be
\emph{stable over time} so that past observations predict future ones.

\textbf{Strict stationarity:} The joint distribution of $\{X_{t_1},\ldots,X_{t_k}\}$
is identical to that of $\{X_{t_1+l},\ldots,X_{t_k+l}\}$ for all $l$. Rarely
verified in practice.

\textbf{Weak (2nd-order) stationarity:} The practical version. $\{X_t\}$ is weakly
stationary if both:
\begin{itemize}
    \item The mean $E(X_t) = \mu$ is constant (does not depend on $t$).
    \item The covariance $\text{Cov}(X_t, X_{t-l}) = v_l$ depends only on the
      \emph{lag} $l$, not on $t$.
\end{itemize}

\textbf{Examples of non-stationary processes:}
\begin{itemize}
    \item \textbf{Structural break:} the mean shifts at some time $k$
      (e.g., $y_t = \beta + e_t$ for $t \leq k$, $y_t = \beta + \lambda + e_t$
      for $t > k$).
    \item \textbf{Random walk (unit root):} $y_t = y_{t-1} + e_t$.
      Variance grows without bound: $\text{Var}(y_t) > \text{Var}(y_{t-1})$,
      so it is not stationary.
\end{itemize}

\newpage
%------------------------------------------------------------
\section{Auto-Correlation and the AR Model}
%------------------------------------------------------------

\subsection{From Correlation to Auto-Correlation}
Apply the correlation formula with $X \to X_t$ and $Y \to X_{t-l}$ (the same series
shifted by lag $l$). The result is \textbf{auto-correlation}.

Define the auto-covariance at lag $l$:
\[
v_l = \text{Cov}(X_t,\, X_{t-l}) = E[(X_t - \mu)(X_{t-l} - \mu)]
\]

The \textbf{auto-correlation function (ACF)} is the normalized version:
\[
r_l = \frac{v_l}{v_0}, \quad r_0 = 1,\quad |r_l| \leq 1,\quad r_l = r_{-l}
\]

where $v_0 = \text{Var}(X_t)$.

\textbf{What $r_l$ means:} The correlation $r_l$ between $X_t$ and $X_{t-l}$
quantifies how much knowing $X_{t-l}$ reduces uncertainty about $X_t$, without
any other information. The best linear predictor of $X_t$ from $X_{t-l}$ is:
\[
X_t - \mu \;\approx\; r_l\,(X_{t-l} - \mu)
\]

\textbf{Practical formula for a time series of length $N$:}
\[
\text{ACF}[l] = \frac{1}{N - l}\sum_{k=1}^{N-l}(x[k]-m)(x[k+l]-m)
\]
Note the denominator is $N - l$ (not $N$) because only $N - l$ overlapping
samples exist after shifting by $l$.

\subsection{The Auto-Regressive (AR) Model}
If a stationary time series can be explained as a linear combination of its own past
values plus noise, it is called an \textbf{autoregressive model of order $p$},
written AR($p$):
\[
X_t = a_0 + a_1 X_{t-1} + a_2 X_{t-2} + \cdots + a_p X_{t-p} + \varepsilon_t
\]
where $\varepsilon_t$ is a zero-mean random noise term and $a_i$ are coefficients.

\subsection{Properties of the AR(1) Model}
The simplest case, AR(1): $X_t = a_0 + a_1 X_{t-1} + \varepsilon_t$.

\begin{itemize}
    \item \textbf{Mean:} $E(X_t) = \mu = \dfrac{a_0}{1 - a_1}$
      (exists only when $a_1 \neq 1$).
    \item \textbf{Variance:} $\text{Var}(X_t) = v_0 = \dfrac{\sigma_\varepsilon^2}{1 - a_1^2}$
    \item \textbf{Stationarity condition:} $|a_1| < 1$.
      If $|a_1| \geq 1$, the variance is infinite or negative --- the series
      explodes or is undefined.
    \item \textbf{ACF decay:} $r_l = a_1^l$. The autocorrelation of a stationary
      AR(1) series decays \textbf{exponentially} with rate $a_1$.
\end{itemize}

\textit{Intuition:} If $a_1 = 0.9$, yesterday's value explains 90\% of today's
fluctuation around the mean. If $a_1 = 0.1$, the series is nearly memoryless
--- past values barely matter.

\subsection{Properties of the AR(2) Model}
AR(2): $X_t = a_0 + a_1 X_{t-1} + a_2 X_{t-2} + \varepsilon_t$.

\begin{itemize}
    \item \textbf{Mean:} $\mu = \dfrac{a_0}{1 - a_1 - a_2}$
      (requires $a_1 + a_2 \neq 1$).
    \item \textbf{ACF recursion:} $r_l = a_1 r_{l-1} + a_2 r_{l-2}$ for $l \geq 2$.
    \item \textbf{Lag-1 ACF:} $r_1 = \dfrac{a_1}{1 - a_2}$
    \item The ACF of a stationary AR(2) satisfies the second-order difference equation
      $(1 - a_1 B - a_2 B^2)\,r_l = 0$, where $B$ is the back-shift operator
      ($B r_l = r_{l-1}$).
\end{itemize}

\subsection{The Lag Operator and ARMA (Reference)}
The \textbf{lag operator} $L$ is defined by $L X_t = X_{t-1}$. Using it, AR($p$)
becomes $a(L) X_t = \varepsilon_t$ where $a(L) = 1 - a_1 L - \cdots - a_p L^p$.

An \textbf{ARMA($p,q$)} model combines autoregressive and moving-average terms:
\[
a(L)\,y_t = b(L)\,e_t
\]
where $b(L) = 1 + \theta_1 L + \cdots + \theta_q L^q$.
This framework unifies AR and MA models and forms the backbone of classical
time series forecasting.

\newpage
%------------------------------------------------------------
\section{Missing Data: A Preview for Next Session}
%------------------------------------------------------------

The transcript closed with a preview of the \emph{missing data} discussion coming
next week. The key point is that the right answer is always \textbf{``it depends''},
not a single prescribed action. The questions to ask are:

\begin{enumerate}
    \item \textbf{What fraction is missing?} --- 0.1\% vs.\ 30\% call for very
      different treatments.
    \item \textbf{Which variables are missing?} --- Categorical or continuous?
      High-importance or low-importance features?
    \item \textbf{How are they missing?} --- Missing in a contiguous range? Missing
      correlated with another variable? Independently at random?
\end{enumerate}

Giving this structured answer in an interview signals that you understand data
qualitatively, not just algorithmically. XGBoost handles missing values natively
(Section~\ref{sec:why_xgb_performs}), but understanding \emph{why} the data is
missing is still the analyst's job.

\newpage
%------------------------------------------------------------
\section{Key Formulas Reference}
%------------------------------------------------------------

\begin{description}
    \item[Pearson Correlation:]
    $r_{x,y} = \dfrac{\text{Cov}(X,Y)}{\sqrt{\text{Var}(X)\,\text{Var}(Y)}}$,
    \quad $-1 \leq r \leq 1$

    \item[Auto-Covariance (lag $l$, length $N$):]
    $\text{ACF}[l] = \dfrac{1}{N-l}\displaystyle\sum_{k=1}^{N-l}(x[k]-m)(x[k+l]-m)$

    \item[Auto-Correlation:]
    $r[l] = \dfrac{\text{ACF}[l]}{\text{ACF}[0]}$,
    \quad $r[0]=1$, \quad $|r[l]|\leq 1$

    \item[AR($p$) Model:]
    $X_t = a_0 + a_1 X_{t-1} + \cdots + a_p X_{t-p} + \varepsilon_t$

    \item[AR(1) Mean:]
    $\mu = \dfrac{a_0}{1-a_1}$

    \item[AR(1) Variance:]
    $\text{Var}(X_t) = \dfrac{\sigma_\varepsilon^2}{1-a_1^2}$, \quad requires $|a_1|<1$

    \item[AR(1) ACF Decay:]
    $r_l = a_1^l$ \quad (exponential decay)

    \item[XGBoost Objective:]
    $\mathcal{L}(\phi) = \displaystyle\sum_i l(\hat{y}_i, y_i) + \sum_k \Omega(f_k)$,
    \quad $\Omega(f) = \gamma T + \tfrac{1}{2}\lambda\|w\|^2$

    \item[Gradient Tree Boosting (2nd-order approximation):]
    $\mathcal{L}^{(t)} \approx \displaystyle\sum_{i=1}^n
    \!\left[g_i f_t(\mathbf{x}_i) + \tfrac{1}{2}h_i f_t^2(\mathbf{x}_i)\right]
    + \Omega(f_t)$
    where $g_i = \partial_{\hat{y}^{(t-1)}} l(y_i, \hat{y}^{(t-1)})$
    and $h_i = \partial^2_{\hat{y}^{(t-1)}} l(y_i, \hat{y}^{(t-1)})$
\end{description}

\end{document}
